\section{Historia académica marginal de la inteligencia corporizada}

La sección anterior explicó los límites entre el lenguaje y la visión; esta sección entra en la historia académica. Wang He (王鹤) subraya que «inteligencia corporizada» y «robótica» no fueron, en su origen, un mismo relato sostenido por un mismo grupo de investigadores. La inteligencia corporizada surgió sobre todo de la búsqueda, por parte de los investigadores de visión por computadora, del problema de la siguiente etapa: pasar de la passive perception a agentes capaces de actuar, de cambiar lo que observan y de aprender en lazo cerrado. Los investigadores de la tradición robótica, en cambio, se han ocupado durante mucho tiempo del control, la estructura mecánica, la optimización, los escenarios industriales y el hardware concreto.

Esta diferencia explica por qué la tarea más importante en los inicios de la inteligencia corporizada no fue la manipulación diestra, sino la navegación. La navegación permite a los investigadores de visión permanecer en una zona de confort relativamente conocida: el agent cambia su propia posición, la cámara capta una imagen nueva y el entorno no tiene que modificarse. Ya cuenta con un ciclo de percepción-acción, pero es más fácil de estandarizar que el agarre, la apertura de puertas, el contacto y el control de fuerza. Por ello, tareas y plataformas como Object Goal Navigation, Point Goal Navigation y Habitat se convirtieron en el centro del relato inicial.

\begin{figure}[H]
\centering
\includegraphics[width=0.96\textwidth]{embodied-ai-origin.png}
\caption{Historia académica marginal de la inteligencia corporizada: de la visión, la navegación y el ciclo de percepción y acción hacia la corriente principal. Diagrama conceptual propio, elaborado a partir de la conversación entre 00:25:08 y 00:41:15.}
\end{figure}

\begin{knowledgebox}{Lectura de la figura: la corriente académica no surgió de manera natural de la «industria robótica»}
En la figura, el recorrido de Vision a Navigation y después a Perception y Action muestra que la inteligencia corporizada fue, en su origen, una extensión de la «siguiente generación de tareas de visión» propuesta por investigadores de visión. Más adelante incorporó robot learning, el agarre, el aprendizaje por refuerzo y los modelos grandes, pero su relato central fue al principio este: la visión no puede limitarse a reconocer de forma pasiva, sino que debe integrarse en un lazo cerrado de acción.
\end{knowledgebox}

\subsection{Perception-Action Loop}

Esta subsección expone con claridad el paradigma de fondo de la inteligencia corporizada. El Perception-Action Loop consiste en «decidir primero la acción a partir de la percepción; la acción modifica el entorno o la propia posición, y la nueva percepción guía el paso siguiente». La navegación es el ejemplo más sencillo: el robot avanza, cambian sus coordenadas y el ángulo de la cámara, y en consecuencia ve un entorno nuevo. Las tareas de manipulación son más complejas: después de tomar una botella, cambian la posición de la botella, el estado de lo que se sostiene en la mano y las acciones posibles a continuación.

En comparación con el agarre de un solo paso, el lazo cerrado refleja mejor el valor de investigación de la inteligencia corporizada. Muchos de los primeros métodos de agarre convertían el problema en razonamiento sobre nubes de puntos o geometría: observar el objeto una vez, predecir el punto de agarre, ejecutar una sola vez y dar la tarea por terminada. Aunque esto ocurre en un robot, sigue pareciéndose a una tarea de reconocimiento de patrones tridimensionales. Lo que persigue la inteligencia corporizada es el lazo cerrado de varios pasos: cada acción modifica el estado del mundo, y el modelo debe volver a observar, volver a planificar y seguir ejecutando.

\begin{figure}[H]
\centering
\includegraphics[width=0.96\textwidth]{perception-action-loop.png}
\caption{Perception-Action Loop: el relato central de la inteligencia corporizada es el lazo cerrado de percepción, acción y retroalimentación. Diagrama conceptual propio, elaborado a partir de la conversación entre 00:25:08 y 00:41:15.}
\end{figure}

\begin{knowledgebox}{Lectura de la figura: el lazo cerrado convierte la visión de clasificación en acción}
De Observe a Infer, y después a Act, Feedback y Adapt, el cambio decisivo es que el entorno ya no es solo una imagen de entrada, sino que genera estados nuevos como consecuencia de las acciones. Al leer la figura conviene prestar atención al último paso, Adapt: la inteligencia corporizada no es una predicción única, sino la corrección continua de la política mediante la retroalimentación.
\end{knowledgebox}

\begin{importantbox}{Hitos representativos}
Wang He menciona que Li Feifei (李飞飞) calificó la inteligencia corporizada como una de las tres estrellas polares del futuro de la visión por computadora, y que Jensen Huang (黄仁勋), en eventos de NVIDIA, ha insistido en que la próxima generación de IA es la embodied AI. Este tipo de declaraciones llevó una dirección académica marginal hacia la corriente principal e hizo que las comunidades de visión, robótica, aprendizaje por refuerzo y modelos grandes comenzaran a converger en torno a un mismo término.
\end{importantbox}

\subsection{Resumen de la sección}

La historia académica de la inteligencia corporizada es una migración hacia una «visión activa»: del reconocimiento de imágenes a la navegación, del razonamiento de un solo paso a la acción en lazo cerrado, y después la incorporación gradual de la manipulation, el robot learning y la capacidad de los modelos grandes para tareas generales.
